\def\tituloconcepto{Decoherencia y spin echo}
% Si el capitulo se compila solo (sin main.tex), esta cabecera abre el documento.
% Antes de % Si el capitulo se compila solo (sin main.tex), esta cabecera abre el documento.
% Antes de % Si el capitulo se compila solo (sin main.tex), esta cabecera abre el documento.
% Antes de \input{cabecera} cada capitulo define \tituloconcepto.
\ifdefined\enmain\else
  \documentclass[11pt,a4paper]{article}
  \input{preambulo}
  \begin{document}
  \title{\tituloconcepto}
  \author{\autor}
  \date{\today}
  \maketitle
\fi
 cada capitulo define \tituloconcepto.
\ifdefined\enmain\else
  \documentclass[11pt,a4paper]{article}
  % Preambulo compartido por todos los documentos. Compilar con XeLaTeX/LuaLaTeX
% (o Tectonic). Si se usa Overleaf, elegir el compilador XeLaTeX.
\usepackage[spanish,es-noshorthands]{babel}
\usepackage{amsmath,amssymb}
\usepackage{booktabs,array}
\usepackage[a4paper,margin=2.5cm]{geometry}
\usepackage{xcolor}
\usepackage{listings}
\usepackage[most]{tcolorbox}
\usepackage{quantikz}
\usepackage[hidelinks]{hyperref}

% ---- Autor (editar a gusto) -------------------------------------------------
\newcommand{\autor}{Genesis Silesky}

% ---- Notacion de Dirac y atajos ---------------------------------------------
% quantikz ya define \ket y \bra; DeclareRobustCommand los redefine sin error.
\DeclareRobustCommand{\ket}[1]{\left|#1\right\rangle}
\DeclareRobustCommand{\bra}[1]{\left\langle#1\right|}
\DeclareRobustCommand{\braket}[2]{\left\langle#1\middle|#2\right\rangle}
\newcommand{\C}{\mathbb{C}}
\newcommand{\R}{\mathbb{R}}
\DeclareMathOperator{\Tr}{Tr}

% ---- Cajas de apoyo --------------------------------------------------------
% "concepto": definiciones previas que hacen falta para entender el capitulo.
\newtcolorbox{concepto}[1]{breakable,colback=blue!4,colframe=blue!55!black,
  fonttitle=\bfseries,title={Concepto previo: #1}}
% "nota": aclaraciones y advertencias.
\newtcolorbox{nota}[1]{breakable,colback=orange!6,colframe=orange!70!black,
  fonttitle=\bfseries,title={#1}}

% ---- Estilos de listados ----------------------------------------------------
\lstdefinestyle{py}{language=Python,basicstyle=\ttfamily\footnotesize,
  breaklines=true,postbreak=\mbox{\textcolor{gray}{$\hookrightarrow$}\space},
  numbers=left,numberstyle=\tiny\color{gray},frame=single,
  commentstyle=\color{green!40!black}\itshape,
  keywordstyle=\color{blue!70!black}\bfseries,
  stringstyle=\color{red!60!black},showstringspaces=false,tabsize=4,
  columns=fullflexible,keepspaces=true}
\lstdefinestyle{salida}{basicstyle=\ttfamily\footnotesize,breaklines=true,
  frame=leftline,backgroundcolor=\color{gray!8},columns=fullflexible,
  keepspaces=true}

  \begin{document}
  \title{\tituloconcepto}
  \author{\autor}
  \date{\today}
  \maketitle
\fi
 cada capitulo define \tituloconcepto.
\ifdefined\enmain\else
  \documentclass[11pt,a4paper]{article}
  % Preambulo compartido por todos los documentos. Compilar con XeLaTeX/LuaLaTeX
% (o Tectonic). Si se usa Overleaf, elegir el compilador XeLaTeX.
\usepackage[spanish,es-noshorthands]{babel}
\usepackage{amsmath,amssymb}
\usepackage{booktabs,array}
\usepackage[a4paper,margin=2.5cm]{geometry}
\usepackage{xcolor}
\usepackage{listings}
\usepackage[most]{tcolorbox}
\usepackage{quantikz}
\usepackage[hidelinks]{hyperref}

% ---- Autor (editar a gusto) -------------------------------------------------
\newcommand{\autor}{Genesis Silesky}

% ---- Notacion de Dirac y atajos ---------------------------------------------
% quantikz ya define \ket y \bra; DeclareRobustCommand los redefine sin error.
\DeclareRobustCommand{\ket}[1]{\left|#1\right\rangle}
\DeclareRobustCommand{\bra}[1]{\left\langle#1\right|}
\DeclareRobustCommand{\braket}[2]{\left\langle#1\middle|#2\right\rangle}
\newcommand{\C}{\mathbb{C}}
\newcommand{\R}{\mathbb{R}}
\DeclareMathOperator{\Tr}{Tr}

% ---- Cajas de apoyo --------------------------------------------------------
% "concepto": definiciones previas que hacen falta para entender el capitulo.
\newtcolorbox{concepto}[1]{breakable,colback=blue!4,colframe=blue!55!black,
  fonttitle=\bfseries,title={Concepto previo: #1}}
% "nota": aclaraciones y advertencias.
\newtcolorbox{nota}[1]{breakable,colback=orange!6,colframe=orange!70!black,
  fonttitle=\bfseries,title={#1}}

% ---- Estilos de listados ----------------------------------------------------
\lstdefinestyle{py}{language=Python,basicstyle=\ttfamily\footnotesize,
  breaklines=true,postbreak=\mbox{\textcolor{gray}{$\hookrightarrow$}\space},
  numbers=left,numberstyle=\tiny\color{gray},frame=single,
  commentstyle=\color{green!40!black}\itshape,
  keywordstyle=\color{blue!70!black}\bfseries,
  stringstyle=\color{red!60!black},showstringspaces=false,tabsize=4,
  columns=fullflexible,keepspaces=true}
\lstdefinestyle{salida}{basicstyle=\ttfamily\footnotesize,breaklines=true,
  frame=leftline,backgroundcolor=\color{gray!8},columns=fullflexible,
  keepspaces=true}

  \begin{document}
  \title{\tituloconcepto}
  \author{\autor}
  \date{\today}
  \maketitle
\fi


\section{Decoherencia y spin echo: perder y recuperar la fase de un qubit}

\subsection{Objetivo}
Entender por qué un qubit pierde información con el tiempo (relajación $T_1$ y
$T_2$, estados mezclados), cómo el \emph{spin echo} recupera la fase perdida por
la inhomogeneidad del campo y para qué sirve la secuencia CPMG.

\subsection{Conceptos previos}
\begin{concepto}{esfera de Bloch, precesión y estados mezclados}
Ver capítulos~10 (esfera de Bloch), 11 (precesión y pulsos) y~12 (rotaciones).
Un \emph{estado puro} es un punto en la superficie de la esfera de Bloch; un
\emph{estado mezclado} es un punto en su interior.
\end{concepto}

\subsection{Fundamento teórico}
\paragraph{Relajación y pérdida de información.} Un qubit de espín $1/2$ en
un campo magnético tiene $\ket\uparrow\leftrightarrow\ket0$ y
$\ket\downarrow\leftrightarrow\ket1$. Idealmente, el módulo del momento
angular de espín promedio, $|\langle S\rangle|=\sqrt{\langle S_x\rangle^2+\langle
S_y\rangle^2+\langle S_z\rangle^2}$, se conserva y el estado es puro. En la
práctica, la relajación hace que el vector se acorte y entre en la esfera:
el estado se vuelve \textbf{mezclado},
$\rho_{\mathrm{mix}}=p_1\ket1\bra1+p_2\ket2\bra2+\cdots$, una combinación de
estados puros con probabilidades clásicas. Eso representa \emph{pérdida de
información}, porque aumenta la incertidumbre. Si el estado se mezcló poco,
se puede estimar el estado puro correspondiente; si la componente transversal
desaparece por completo, la información de fase se perdió por completo, y en el
caso extremo (estado máximamente mezclado) ya no se puede obtener ninguna
información midiendo. Hay dos procesos:
\begin{itemize}
  \item \textbf{Relajación longitudinal} ($T_1$): cambia la población de los
    niveles (la componente $z$). Es la pérdida de energía hacia el entorno.
  \item \textbf{Relajación transversal} ($T_2$, \emph{decoherencia}): se pierde la
    coherencia de fase (las componentes $x$ e $y$).
\end{itemize}
En los sistemas de resonancia magnética nuclear (RMN), la relajación transversal
pesa más que la longitudinal en la pérdida de información.

\paragraph{Inhomogeneidad del campo.} Una muestra tiene muchos espines (un
\emph{conjunto}). Si el campo no es uniforme, cada espín gira con una
frecuencia de Larmor ligeramente distinta, así que sus fases se separan con el
tiempo y la magnetización neta (el promedio) se acorta. Tras un pulso
$R_y(\pi/2)$ que lleva el espín a $x$, la señal oscila y decae con envolvente
exponencial $e^{-t/T_2^{*}}$. En un sistema sin conjunto (un solo espín medido
muchas veces) pasa lo mismo si el campo fluctúa en el \emph{tiempo}.

\paragraph{El eco de espín.} Erwin Hahn descubrió en 1950 que un pulso $\pi$ (de 180
grados) aplicado en medio de la relajación revierte las fases y produce un
``eco''. Para un conjunto de espines, con tiempo total $t$:
\begin{enumerate}
  \item Al inicio todos apuntan a $x$: polarización máxima.
  \item Durante $t/2$ cada espín acumula la fase $\theta_i=\omega_it/2$ y la
    magnetización decae.
  \item Un pulso $R_x(\pi)$ en $t/2$ invierte las fases: $\theta_i\to-\omega_it/2$.
  \item Durante la segunda mitad cada espín vuelve a acumular $\omega_it/2$.
  \item En $t$ la fase es $-\omega_it/2+\omega_it/2=0$ para todos: la polarización
    se recupera.
\end{enumerate}
El eco es \emph{menor} que la señal inicial: el spin echo retrasa la
relajación transversal pero no la elimina, porque los espines también
interactúan con la red (relajación natural). Con eco, la señal decae con
$T_2$; sin eco, con $T_2^{*}$, y en general $T_2\ge T_2^{*}$. Después de
$t=T_2$ la componente transversal baja a $1/e\approx0.368$ de su máximo. Para
informática cuántica se usa $T_2$ como tiempo de coherencia.

\paragraph{Secuencias.} Un solo pulso $\pi$ da un eco, pero su fase puede no
coincidir con la inicial (depende del eje del pulso y de la dirección inicial).
Para conservar el estado se añade un segundo pulso $\pi$ en el eco: así la fase
coincide siempre con la inicial. La secuencia \textbf{CPMG} (Carr, Purcell,
Meiboom y Gill) aplica pulsos $\pi$ a intervalos regulares para intervalos
largos; \emph{CPMG-$n$} son $n$ pulsos en un tiempo $t$, con separación
$t/n$ (el primero y el último a $t/2n$ de los extremos). CPMG funciona en muchos
experimentos, pero no elimina todas las fuentes de decaimiento (como la
difusión o el intercambio químico).

\begin{nota}{Fuentes y modelo del ejemplo}
Este capítulo se basa en las notas de spin echo del curso. Los valores
numéricos ($T_1=100$~ms, $T_2=40$~ms, dispersión de frecuencias de $100$~Hz,
fluctuación lenta de $30$~Hz con tiempo de correlación de $10$~ms) \textbf{son
elegidos para ilustrar}, no medidos. El archivo de teoría de decoherencia del
curso es un escaneo sin texto legible, así que el modelo de Bloch para $T_1$ y
$T_2$ es el estándar de la definición de las notas.
\end{nota}

\subsection{Código (completamente comentado)}
\lstinputlisting[style=py]{../../fuentes/21_decoherencia_spin_echo.py}

\subsection{Explicación celda por celda}
\begin{description}
  \item[Celda 2] Evoluciona el vector de Bloch con $T_1$ y $T_2$ y calcula su
    longitud y la pureza $(1+|\mathbf r|^2)/2$.
  \item[Celda 3] Simula 5000 espines con frecuencias distintas y compara la
    magnetización con la fórmula gaussiana $e^{-(2\pi\sigma t)^2/2}$.
  \item[Celda 4] Compara la señal sin eco y con eco para varios tiempos.
  \item[Celda 5] Sigue la fase de un espín con rotaciones $3\times3$ para
    cuatro combinaciones de fase inicial y eje del pulso $\pi$, con uno y
    con dos pulsos.
  \item[Celda 6] Compara sin pulsos, spin echo y CPMG-8 con un campo que
    fluctúa lentamente.
\end{description}

\subsection{Salida obtenida}
\lstinputlisting[style=salida]{../../fuentes/salidas/21_decoherencia_spin_echo.txt}
Usa semilla fija. Con otra semilla cambian ligeramente los números del
conjunto de espines.

\subsection{Interpretación}
La componente transversal baja a $0.368$ en $t=T_2=40$~ms, mientras que la
poblacional ($z$) se recupera más despacio. El desfase de un conjunto coincide con
la fórmula (por ejemplo $0.459$ frente a $0.454$ a los 2~ms). Con eco, la señal
a 32~ms es $0.449$ (solo la relajación natural $e^{-t/T_2}$), mientras que
sin eco es $0.009$. Con un solo pulso, la fase del eco coincide con la inicial solo en 2 de
los 4 casos ($0^\circ$, $180^\circ$, $270^\circ$ y $90^\circ$ para fases
iniciales $0^\circ$ y $90^\circ$ con $\pi_x$ y $\pi_y$); con el segundo pulso
del mismo eje es siempre la inicial. Con un campo que fluctúa lentamente, un solo spin echo no
basta ($0.009$), pero CPMG-8 conserva $0.749$ de la señal.

\subsection{Ejercicios}
\begin{enumerate}
  \item Cambia $T_2$ a $20$~ms y observa cómo cambia la columna ``con eco''.
  \item Prueba CPMG con $n=2,4,16$ y compara la señal final.
  \item ¿Qué pasa con CPMG si el ruido fluctúa muy rápido (tiempo de
    correlación de $1$~ms)?
\end{enumerate}

\subsection{Errores comunes}
\begin{itemize}
  \item Creer que el eco recupera \emph{toda} la señal: solo corrige el
    desfase por campos estáticos (o lentos); la relajación natural $T_2$ sigue.
  \item Confundir $T_2$ con $T_2^{*}$: $T_2^{*}$ incluye la inhomogeneidad y es
    más corto.
\end{itemize}

\ifdefined\enmain\else
  \end{document}
\fi

